\section{Threats to Validity}
\label{sec:threats_validity}

To ensure sound scientific interpretation, we examine the threats to validity that may influence our empirical conclusions.

\subsection{Internal Validity}
Internal validity concerns whether the observed performance differences can be strictly attributed to the mechanisms under evaluation rather than confounding factors.
\begin{itemize}
    \item \textbf{Controlled Ablations:} Each ablation variant was implemented using isolated evaluation adapters (\texttt{DisabledContradictionDetector}, \texttt{NoConfidenceMemoryClassifier}, tier-neutral scoring weights). The core SQLite repository, schema, and baseline retrieval were held identical across all conditions.
    \item \textbf{Deterministic Model Adapters:} To eliminate stochastic inference jitter and API drift, experiments used deterministic mock and rule adapters with a fixed random seed (42). This guarantees that every method processed identical inputs under identical conditions.
    \item \textbf{Risk of Benchmark Overfitting:} The evaluation scenarios were designed to probe specific structural failure modes (contradictions, ambiguous confidence, verbatim constraints). While this allows clean causal isolation, it may not reflect the full variance of unconstrained human dialogues.
\end{itemize}

\subsection{Construct Validity}
Construct validity assesses whether our quantitative metrics accurately measure the theoretical properties of interest.
\begin{itemize}
    \item \textbf{Fact Recall & Answer Correctness:} We evaluated fact recall by checking for the presence of ground-truth target facts within compiled context $C_t$, and answer correctness by substring and semantic matching against target outputs. In rare edge cases, a model might produce a correct answer through parametric pre-training rather than retrieved context; however, our synthetic entities and constraints were designed to minimize this leakage.
    \item \textbf{Token Counting:} Tokens were measured using a deterministic whitespace and punctuation tokenizer. While token counts align closely with standard subword tokenizers (e.g., cl100k\_base), small absolute token offsets may occur across different model families.
\end{itemize}

\subsection{External Validity}
External validity concerns the generalizability of our findings to real-world applications.
\begin{itemize}
    \item \textbf{Synthetic Micro-Benchmark:} Our benchmark comprises 11 controlled scenarios and 22 queries evaluated across 5 budgets (880 trials). While larger than many prototype studies, it does not encompass the full linguistic entropy and conversational noise of web-scale datasets (such as the 35-session dialogues in LoCoMo).
    \item \textbf{Lexical Retrieval vs. Dense Embeddings:} ManiGit's current reference implementation utilizes Okapi BM25 and lexical term matching. While this guarantees deterministic, reproducible retrieval without external vector database dependencies, real-world systems often benefit from dense semantic vector embeddings or hybrid retrieval.
    \item \textbf{Absence of Human Evaluation:} All answers were graded programmatically against ground truth. Blinded human rater evaluations have not yet been conducted.
\end{itemize}

\subsection{Statistical Validity}
Statistical validity addresses the appropriateness of our statistical inference.
\begin{itemize}
    \item \textbf{Unit of Analysis:} Our 880 trials originate from 22 queries evaluated across 8 methods and 5 budgets. We explicitly treat the $N=110$ query-budget combinations as the unit of evaluation and do not falsely claim 110 independent conversations.
    \item \textbf{Confidence Intervals:} We compute exact 95\% Wilson score confidence intervals for all binary metrics, providing rigorous uncertainty bounds appropriate for finite sample sizes without assuming Gaussian normality.
\end{itemize}
